\documentstyle[%


righttag%


,11pt%


,amssymb%


,russian%               remove in Eng.


,izvru%                 Set izven% in Eng.


% ,izvru-ks             for brief communications


]{amsart}





%


% \adjust \hspace, lines 168, 191


% Search \times and join in eng.


\newcommand{\myfont}{\series{m}\shape{n}\selectfont}


\newcommand{\re}{\operatorname{Re}}


\newcommand{\sn}{\operatorname{sn}}


\newcommand{\cn}{\operatorname{cn}}


\newcommand{\im}{\operatorname{Im}}


%\newcommand{\ov}{\overline}


\newcommand{\tts}[1]{}





\theoremstyle{plain}


\theorembodyfont{\it}


\newtheorem{thmnonum}{\hskip\parindent Теорема}


\renewcommand{\thethmnonum}{}





%\hoffset=-15mm%                  For Eng.





\setcounter{page}{35}


\god{1998}


\nomer{5~(432)} %                        Remove in Eng.


%\nomer{Vol.\,42, No.\,5}%               Set in Eng.


%\str{\pageref{begin}--\pageref{end}}%  Set in Eng.


\udk{517.518}


\author{\it А.Л.\,Лукашов}


% NB:   ^^ remove in Eng.


\title{Рациональные интерполяционные процессы\\


на двух отрезках}





\date{}


% \pagestyle{myheadings}%         Set in Eng.





\pagestyle{plain} %              Remove in Eng.


\begin{document}





% \thispagestyle{plain}%          Set in Eng.





\maketitle


\label{begin}





Пусть $E$ --- компактное множество действительных чисел,


$\frak M=\{x_{k,n}\}^{n,\infty}_{k=1,\,n=1}$ --- матрица


узлов интерполирования, лежащих в $E$, $f$ --- произвольная непрерывная


на $E$ действительнозначная функция. Интерполяционный процесс Лагранжа


$\{\cal L_n(\frak M,f,x)\}^{\infty}_{n=1}$ функции $f$


задается формулами


\begin{equation}


\cal L_n(\frak M,f,x)=\sum^n_{k=1}f(x_{k,n})l_{k,n}


(\frak M,x),\quad n=1,2,\dotsc,


\end{equation}


где $l_{k,n}(\frak M,x)$ --- фундаментальные полиномы Лагранжа,


\begin{gather}


l_{k,n}(\frak M,x)=


\frac{\omega_n(x)}{(x-x_{k,n})\omega'_n(x_{k,n})},\quad


k=\ov{1,n},\\


\omega_n(x_{k,n})=\prod^n_{k=1}(x-x_{k,n}),\quad


n=1,2,\dotsc


\end{gather}





Хорошо известна (напр., \cite{1}, \cite{2}) роль функций Лебега


$$


\lambda_n(\frak M,x)=\sum^n_{k=1}|l_{k,n}(\frak M,x)|


$$


и констант Лебега


$$


\lambda_n(\frak M)=\|\lambda_n(\frak M,\cdot)\|_{C(E)}


$$


в изучении поведения интерполяционных процессов Лагранжа. При этом к


наиболее исследованным относится случай $E=[-1,\,1]$ и


$\omega_n(x)=T_n(x)=\cos(n\arccos x)$, т.\,е. матрицы узлов Чебышева.





Для интерполирования рациональными функциями с фиксированными полюсами,


заданными матрицей обратных величин полюсов


$\frak A=\{a_{\lambda,n}\}^{n,\infty}_{\lambda=1,\,n=1}$,


интерполяционный


процесс Лагранжа задается формулами (1)--(2), в которых вместо (3)


нужно положить


$$


\omega_n(x)=\prod\limits^n_{k=1}(x-x_{k,n})/(1-a_{k,n}x).


$$





В комплексном случае полученный процесс, который будет далее


обозначаться через\linebreak


$\{\cal L_n(\frak M,\frak A,f,x)\}^\infty_{n=1}$, достаточно хорошо


исследован для различных классов аналитических функций $f$ (напр.,


\cite{3}, \cite{4}), в то


же время эта теория для действительного случая разработана значительно


меньше.





В частности, для $E=[-1,\,1]$ естественным аналогом узлов Чебышева


являются узлы


Чебышева-Маркова, т.\,е.


$\omega_n(x)=M_n(x)=\cos\Bigl(\sum\limits^n_{k=1}\arccos


\frac{x-a_{k,n}}{1-a_{k,n}x}\Bigr)$. Такие процессы


были введены В.Н.\,Русаком


\cite{5}, и были использованы им, Е.А.\,Ровбой \cite{6}, \cite{7},


А.П.\,Старовойтовым


\cite{8} для интерполирования различных классов


функций. В частности, в работе А.П.\,Старовойтова \cite{8} найдена


оценка соответствующих констант Лебега (многие приемы из этой работы


будут использованы ниже при доказательстве теоремы).





Цель данной работы --- найти оценку констант Лебега


$\lambda_n(\frak M,\frak A)$ для $E=[-1,a]\cup[b,1]$,


$-1<a<b<1$,


матрицы $\frak M$, задаваемой нулями рациональной функции


\begin{equation}


M_n(x)=\frac{P_n(x)}{\prod\limits^n_{k=1}(1-a_{k,n}x)},


\end{equation}


наименее уклоняющейся от нуля на $E$ в равномерной метрике среди всех


дробей вида (4), где


$P_n(x)=x^n+b_1x^{n-1}+\dotsb+b_n$, $b_i\in\Bbb R$, $i=\ov{1,n}$, а


$a_{k,n}$ фиксированы.


При этом потребовалось разработать специальную технику использования


геометрических характеристик величин, выражающихся в эллиптических


функциях.





Будем считать, что величины $a_{k,n}$, $k=1,\dotsc,n$, занумерованы


следующим образом:


$a_{1,n}=\dotsb=a_{\varkappa_n,n}=0$; $a_{i,n}\ne 0$,


$i=\varkappa_n+1,\dotsc,n$; $a^{-1}_{j,n}\in\Bbb C\setminus[-1,1]$,


$j=\varkappa_n+1,\dotsc,n$; $\re a_{j,n}\ge 0$, $j=1,\dotsc,n_1$;


$\re a_{j,n}<0$, $j=n_1+1,\dotsc,n$


(возможную зависимость $n_1$ от $n$


указывать не станем), причем из $\im a_{j,n}>0$ следует


$a_{j+1,n}=\ov a_{j,n}$, $j=\varkappa_n+1,\dotsc,n$


($1\le\varkappa_n\le n_1\le n$).


Через $C_1,C_2,\dotsc$


будем обозначать положительные постоянные, не зависящие от $n$.





Далее, пусть $K=K(k)$ --- полный эллиптический интеграл 1-го рода,


соответствующий модулю $k=\sqrt{\frac{2(b-a)}{(1-a)(1+b)}}$. Назовем


матрицу $\frak A$


регулярной относительно $E$, если величины $\varphi_n(\frak A,k,a)$,


задаваемые соотношениями


\begin{equation}


\varphi_n(\frak A,k,a)=\frac{1}{K}


\int^{\sqrt{(1-a)/2}}_0\sum^n_{\lambda=1}


\frac{\sqrt{(1+a_{\lambda,n})(1-a a_{\lambda,n})}}


{\sqrt{(1-a a_{\lambda,n}-(1+a_{\lambda,n})u^2)


(1-a a_{\lambda,n}-k^2(1+a_{\lambda,n})u^2)}}du,


\end{equation}


принимают целочисленные значения для всех $n\ge 2$.


Нетрудно видеть, что для


любых заданных величин $a_{1,n},\dotsc,a_{n-1,n}$ можно найти такие


$a_{n,n}$, $n=1,2,\dotsc,$ что матрица $\frak A$ будет регулярной


относительно $E$.





Нам потребуется в дальнейшем следующее представление рациональных


функций (4) для регулярных относительно $E$ матриц $\frak A$.


\medskip





{\bf Лемма 1} (\cite{9}). {\it


Пусть


$\varphi_n(\frak A,k,a)=m\in\Bbb N$. Тогда рациональная функция


$M_n(x)$, наименее уклоняющаяся от нуля на $E$ среди всех функций вида


$(4)$, может быть представлена в одном из следующих видов\,{\em :}





$1.\hspace{3.4cm}\displaystyle


M_n(x)=\frac{M_n}{2}\biggl(\prod^n_{\lambda=1}


\frac{H(u-\rho_{\lambda,n})}{H(u+\rho_{\lambda,n})}+


\prod^n_{\lambda=1}


\frac{H(u+\rho_{\lambda,n})}{H(u-\rho_{\lambda,n})}\biggr)$,





\noindent где


\begin{gather*}


x=\frac{\sn^2u\cdot \cn^2\rho_{0,n}+\cn^2u\cdot \sn^2\rho_{0,n}}


{\sn^2u-\sn^2\rho_{0,n}};\\


\sn^2\rho_{\lambda,n}=


\frac{(1-a)(1+a_{\lambda,n})}{2(1-a a_{\lambda,n})},\quad


\lambda=0,1,\dotsc,n;\quad -K'<\im\rho_{\lambda,n}<K',\quad


\lambda=1,2,\dotsc,n;\\


-K<\re\rho_{\lambda,n}<0,\quad \lambda=0,1,\dotsc,n;\quad


K'=K(k'),\quad k'=\sqrt{1-k^2};\quad


a_{0,n}=0;\quad \varepsilon=\pm 1;\\


M_n=\frac{2\varepsilon}{\prod\limits^n_{\lambda=\varkappa_n+1}


a_{\lambda,n}}\prod^n_{\lambda=\varkappa_n+1}


\frac{H(\rho_{0,n}+\rho_{\lambda,n})}{H(\rho_{0,n}-\rho_{\lambda,n})}\cdot


\biggl(\frac{\Theta^2(0)\Theta^2_1(0)}{2\Theta^2_1(\rho_{0,n})


\Theta^2(\rho_{0,n})}\biggr)^{\varkappa_n};


\end{gather*}





$2.\hspace{1.5cm}\displaystyle


M_n(x)=M_n\cos\int^x_{-1}\sum^n_{\lambda=1}


\frac{(x-c_{\lambda,n})\sqrt{(1-a^2_{\lambda,n})


(1-a_{\lambda,n}a)(1-a_{\lambda,n}b)}}


{(1-a_{\lambda,n}x)(1-a_{\lambda,n}c_{\lambda,n})


\sqrt{(1-x^2)(x-a)(x-b)}}dx$,





\noindent где


$$


c_{\lambda,n}=


\frac{\sqrt{1-a_{\lambda,n}a}\,\Theta'(\rho_{\lambda,n})


\sqrt{(1-a)(1+b)}-a\Theta(\rho_{\lambda,n})


\sqrt{(1-a^2_{\lambda,n})(1-a_{\lambda,n}b)}}


{a_{\lambda,n}\sqrt{1-a_{\lambda,n}a}\,\Theta'(\rho_{\lambda,n})


\sqrt{(1-a)(1+b)}-\Theta(\rho_{\lambda,n})


\sqrt{(1-a^2_{\lambda,n})(1-a_{\lambda,n}b)}},\ \;\lambda=1,\dotsc,n,


$$


причем $a<c_{\lambda,n}<b$ для $\im a_{\lambda,n}=0$, и $c_{\lambda,n}$


лежит на дуге окружности, проходящей через $a$, $b$ и


$a^{-1}_{\lambda,n}$, ограниченной точками $a$ и $b$ и не содержащей


$a^{-1}_{\lambda,n}$, при $\im a_{\lambda,n}>0$ $($в этом случае


$c_{\lambda+1,n}=\ov c_{\lambda,n})$, $\lambda=1,\dotsc,n$.


}





\begin{note}


В последнее время интерес к описанным в лемме 1 дробям, которые,


по-видимому, были известны еще Н.И.\,Ахиезеру (но не выписаны им явно),


возрос из-за открытых связей с ортогональными многочленами и решениями


цепочек Тоды (напр., \cite{10}--\cite{13}).


\end{note}


\addtocounter{note}{-1}





Обозначим


\begin{align*}


\gamma_n(x)&=\sum^n_{\lambda=1}


\frac{(x-c_{\lambda,n})\sqrt{(1-a^2_{\lambda,n})


(1-a_{\lambda,n}a)(1-a_{\lambda,n}b)}}


{(1-a_{\lambda,n}x)(1-a_{\lambda,n}c_{\lambda,n})},\\


\gamma^{(1)}_n(x)&=\sum^{n_1}_{\lambda=1}


\frac{\sqrt{1-|a_{\lambda,n}|}\,|x-c_{\lambda,n}|}


{1-|a_{\lambda,n}|x},\\


\gamma^{(2)}_n(x)&=\sum^n_{\lambda=n_1+1}


\frac{\sqrt{1-|a_{\lambda,n}}|\,|x-c_{\lambda,n}|}


{1-|a_{\lambda,n}|x},\\


S_n(x)&=-\int^x_{-1}\frac{\gamma_n(x)dx}


{\sqrt{(1-x^2)(x-a)(x-b)}},\quad x\in E.


\end{align*}





Сформулируем основной результат работы.





\begin{thmnonum}


Пусть $\frak A\subset\{z:|z|<1\}$ --- регулярная относительно $E$


матрица обратных


величин полюсов, не имеющая на $\{z:|z|=1\}$ других предельных точек,


кроме $\pm 1$,


достигаемых ею по некасательным путям, и удовлетворяющая соотношениям


\begin{gather}


\begin{split}


\sum^{n_1}_{\lambda=1}


\frac{\sqrt{1-|a_{\lambda,n}|}\,t}{1-|a_{\lambda,n}|+t^2}


&\ge C_1\quad\text{при} \quad


\sqrt{1-\max_{1\le\lambda\le n_1}|a_{\lambda,n}|}\le t\le 1; \\


\sum^n_{\lambda=n_1+1}\frac{\sqrt{1-|a_{\lambda,n}|}\,t}


{1-|a_{\lambda,n}|+t^2}&\ge C_2\quad\text{при}\quad


\sqrt{1-\max_{n_1+1\le\lambda\le n}|a_{\lambda,n}|}\le t\le 1;


\end{split}


\\


\min_{x\in E}\frac{\gamma_n(x)}


{\sqrt{(x-a)(x-b)}\,\gamma_n((b+1)/2)}\ge C_3,\quad


n=1,2,\dotsc


\end{gather}


Тогда


$$


\cal L_n(\frak M,\frak A)=O(\ln\|\gamma_n\|_{C(E)}).


$$


\end{thmnonum}





\begin{note}


Условие регулярности естественно для интерполирования на двух отрезках,


т.\,к. обеспечивает принадлежность элементов матрицы $\frak M$


множеству $E$.


\end{note}





\begin{note}


При $a=b$ из теоремы следует упомянутый ранее результат


А.П.\,Старовойтова \cite{8} (условия регулярности и (7) здесь


выполняются автоматически).


\end{note}





\begin{note}


При $a=b$ и $a_{1,n}=\dotsb=a_{n,n}=0$ ((6) здесь также выполнено)


получаем порядковую оценку


констант Лебега интерполяционного процесса Лагранжа по узлам Чебышева,


которая, как показано в \cite{14}, является оптимальной по порядку для


полиномиального интерполирования на отрезке. В случае рационального


интерполирования (т.\,е. процессов


$\{\cal L_n(\frak M,\frak A,f,x)\}^\infty_{n=1}$) аналог этого


результата,


насколько известно автору, отсутствует, хотя в \cite{15} было найдено


характеристическое условие на оптимальную матрицу $\frak M$.


\end{note}


\setcounter{lem}{1}





\begin{lem}


В условиях теоремы найдутся постоянные $C_4$ и $C_5$ такие, что при


всех $n\in\Bbb N$


\begin{equation}


C_4\sum^n_{\lambda=1}


\frac{\sqrt{1-|a_{\lambda,n}}|\,|x-c_{\lambda,n}|}


{1-|a_{\lambda,n}|x}\le|\gamma_n(x)|\le


C_5\sum^n_{\lambda=1}


\frac{\sqrt{1-|a_{\lambda,n}|}\,|x-c_{\lambda,n}|}


{1-|a_{\lambda,n}|x}.


\end{equation}


\end{lem}





\begin{pf}


Нетрудно видеть, что лемма следует из неравенства


\begin{equation}


-\frac{\pi}{2}+\delta<\varphi_\lambda(x)<\frac{\pi}{2}-\delta,


\end{equation}


где $\delta>0$ не зависит от $\lambda$ и $n$, а


$$


\varphi_\lambda(x)=\arg


\frac{\sqrt{(1-a^2_{\lambda,n})(1-a_{\lambda,n}a)


(1-a_{\lambda,n}b)}\,(x-c_{\lambda,n})}


{(1-a_{\lambda,n}x)(1-a_{\lambda,n}c_{\lambda,n})}.


$$





Доказательство неравенства (9) (для $\im a_{\lambda,n}<0$)


разобьем на несколько этапов.





1 этап. Докажем, что


$\arg\frac{x-c_{\lambda,n}}{a^{-1}_{\lambda,n}-x}$ заключен между


$\arg\frac{b-c_{\lambda,n}}{a^{-1}_{\lambda,n}-b}$ и


$\arg\frac{1-c_{\lambda,n}}{a^{-1}_{\lambda,n}-1}$. Для этого применим


отображение $w(z)=\frac{z-c_{\lambda,n}}{a^{-1}_{\lambda,n}-z}$,


переводящее отрезок $[b,1]$ в дугу окружности $L$. Окружность


$L$ пересекает действительную ось в двух точках: $w=-1$ и $w=w(y)$, т.\,е.


$(c_{\lambda,n}-y)/(a^{-1}_{\lambda,n}-y)$


действительное. Последнее соотношение означает (\cite{16}, гл.\,1,


с.\,42), что $y$ лежит на прямой, проходящей через $a^{-1}_{\lambda,n}$


и $c_{\lambda,n}$. Отсюда, из


леммы 1 и того, что $\im a^{-1}_{\lambda,n}>0$ и


$|a^{-1}_{\lambda,n}|>1$, следует $y\in(a,b)$ и


$\frac{y-c_{\lambda,n}}{a^{-1}_{\lambda,n}-y}>0$.


Таким образом,


точка $w=0$ находится внутри окружности $L$, и из определения $w(z)$


следует требуемое.





2 этап. Мы доказали, что $\varphi_\lambda(x)$ находится между


\begin{equation}


\arg\sqrt{\frac{(a^{-2}_{\lambda,n}-1)(a^{-1}_{\lambda,n}-a)}


{a^{-1}_{\lambda,n}-b}}\cdot\frac{b-c_{\lambda,n}}


{a^{-1}_{\lambda,n}-c_{\lambda,n}}


\end{equation}


и


\begin{equation}


\arg\sqrt{\frac{(a^{-1}_{\lambda,n}+1)(a^{-1}_{\lambda,n}-b)}


{a^{-1}_{\lambda,n}-a}}\cdot\frac{1-c_{\lambda,n}}


{a^{-1}_{\lambda,n}-c_{\lambda,n}}.


\end{equation}


Из расположения точек $a$, $b$, $c_{\lambda,n}$ и $a^{-1}_{\lambda,n}$


на одной окружности следует


(\cite{16}, с.\,42), что\linebreak


$\Bigl.\frac{b-c_{\lambda,n}}{a^{-1}_{\lambda,n}-c_{\lambda,n}}


\Bigr/\frac{b-a}{a^{-1}_{\lambda,n}-a}$ является


действительным


числом, причем положительным, как показывают неравенства


$$


0<\arg(a^{-1}_{\lambda,n}-a)<\pi\quad


\text{и}\quad


0<\arg(b-c_{\lambda,n})-\arg(a^{-1}_{\lambda,n}-c_{\lambda,n})<\pi.


$$


Отсюда выражение (10) равно


\begin{equation}


\arg\sqrt{\frac{a^{-2}_{\lambda,n}-1}


{(a^{-1}_{\lambda,n}-a)(a^{-1}_{\lambda,n}-b)}}.


\end{equation}


Аналогично (11) преобразуется к


\begin{equation}


\arg\sqrt{\frac{(a^{-1}_{\lambda,n}+1)


(a^{-1}_{\lambda,n}-a)(a^{-1}_{\lambda,n}-b)}


{a^{-1}_{\lambda,n}-1}}


\cdot\frac{1}{a^{-1}_{\lambda,n}-z},


\end{equation}


где $z$ --- точка пересечения окружности, проходящей через $1$,


$c_{\lambda,n}$ и $a^{-1}_{\lambda,n}$,


с действительной осью, причем $a<z<b$. Следовательно, (13) заключено


между


\begin{equation}


\arg\sqrt{\frac{(a^{-1}_{\lambda,n}+1)


(a^{-1}_{\lambda,n}-a)(a^{-1}_{\lambda,n}-b)}


{a^{-1}_{\lambda,n}-1}}


\cdot\frac{1}{a^{-1}_{\lambda,n}-a}


=\arg\sqrt{\frac{(a^{-1}_{\lambda,n}+1)(a^{-1}_{\lambda,n}-b)}


{(a^{-1}_{\lambda,n}-1)(a^{-1}_{\lambda,n}-a)}}


\end{equation}


и


\begin{equation}


\arg\sqrt{\frac{(a^{-1}_{\lambda,n}+1)


(a^{-1}_{\lambda,n}-a)(a^{-1}_{\lambda,n}-b)}


{a^{-1}_{\lambda,n}-1}}


\cdot\frac{1}{a^{-1}_{\lambda,n}-b}


=\arg\sqrt{\frac{(a^{-1}_{\lambda,n}+1)(a^{-1}_{\lambda,n}-a)}


{(a^{-1}_{\lambda,n}-1)(a^{-1}_{\lambda,n}-b)}}.


\end{equation}


Теперь, оценивая (12), (14) и (15) с помощью вытекающих из стремления


$\{a_{\lambda,n}\}$ к $\pm 1$ по некасательным путям неравенств


\cite{8}


\begin{gather}


C_6(1-|a_{k,n}|x)\le|1-a_{k,n}x|\le C_7(1-|a_{k,n}|x),\\


-\frac{\pi}{2}+\varphi_0\le\arg


\frac{\sqrt{1-a^2_{k,n}}}{1-a_{k,n}x}\le\frac{\pi}{2}-


\varphi_0,\notag


\end{gather}


($\varphi_0>0$ не зависит от $n$), получим (9), а с ним и лемму.


\end{pf}





{\bf Доказательство теоремы.} Будем считать, что величина


$\varphi_n(\frak A,k,a)$ равна $m_n=m$,


тогда $-1<x_{n,n}<\dotsb<x_{m+1,n}<a<b<x_{m,n}<\dotsb<x_{1,n}<1$.


Кроме того, пусть $i$, $p$ и $q$ выбираются из неравенств


$x_{i,n}<x<x_{i-1,n}$, $x_{p+1,n}\le(a-1)/2<x_{p,n}$,


$x_{q+1,n}\le(b+1)/2<x_{q,n}$.





Рассмотрим случай $i=1,\dotsc,q$; $k=i+1,\dotsc,m$, и оценим


\begin{multline}


|l_{k,n}(\frak M,x)|\le


\frac{\sqrt{(1-x^2_{k,n})(x_{k,n}-a)(x_{k,n}-b)}}


{(x-x_{k,n})|\gamma_n(x_{k,n})|}=\\


%


=\frac{|S_n(x_{k-1,n})-S_n(x_{k,n})|


\sqrt{(1-x^2_{k,n})(x_{k,n}-a)(x_{k,n}-b)}}


{(x-x_{k,n})|\gamma_n(x_{k,n})|\pi}.


\end{multline}


Применяя теорему Лагранжа к (17), найдем


\begin{equation}


|l_{k,n}(\frak M,x)|\le\frac{1}{\pi}


\frac{x_{k-1,n}-x_{k,n}}{x-x_{k,n}}\cdot Z_k\cdot H_k,


\end{equation}


где


$$


Z_k=\sqrt{\frac{1-x^2_{k,n}}{1-\xi^2_k}},\quad


H_k=\biggl|\frac{\gamma_n(\xi_k)}{\gamma_n(x_{k,n)}} \biggr|


\frac{\sqrt{(x_{k,n}-a)(x_{k,n}-b)}}


{\sqrt{(\xi_k-a)(\xi_k-b)}},\quad x_{k,n}<\xi_k<x_{k-1,n}.


$$


Для оценки


\begin{equation}


Z_k\le\sqrt{\frac{1-x_{k,n}}{1-\xi_k}}=


\sqrt{1+\frac{\xi_k-x_{k,n}}{1-\xi_k}}


\end{equation}


имеем


\begin{equation}


\frac{\xi_k-x_{k,n}}{1-\xi_k}\le


\frac{x_{k-1,n}-x_{k,n}}{1-\xi_k}\le


\frac{\pi\sqrt{(1-\xi^2_k)(\xi_k-a)(\xi_k-b)}}


{(1-\xi_k)|\gamma_n(\xi_k)|},


\end{equation}


и при $\min((1-b)/2,\,1/2)\ge 1-\xi_k\ge 1-


\max\limits_{1\le\lambda\le n_1}|a_{\lambda,n}|$ правая часть


(20) не превосходит


\begin{equation}


\frac{2\pi\sqrt{(\xi_k-a)(\xi_k-b)}}


{C_4\sqrt{1-\xi_k}\displaystyle\sum\limits^n_{j=1}


\dfrac{|\xi_k-c_{j,n}|\sqrt{1-|a_{j,n}|}}{1-|a_{j,n}|\xi_k}}\le


C_8\biggl(\sum^n_{j=1}\frac{\sqrt{1-|a_{j,n}|}\sqrt{1-\xi_k}}


{1-|a_{j,n}|+1-\xi_k}\biggr)^{-1}\le\frac{C_8}{C_1}.


\end{equation}


При $1-\xi_k>\min((1-b)/2,\,1/2)$ получаем


\begin{equation}


Z_k\le\frac{1}{\sqrt{\min\left(\dfrac{1-b}{2},\dfrac{1}{2}\right)}}.


\end{equation}


Для оценки $Z_k$ при


$0\le 1-\xi_k\le 1-\max\limits_{1\le\lambda\le n_1}|a_{\lambda,n}|$


выберем $\theta_1$ так, чтобы для $1-\theta_1\le x\le 1$ выполнялось


неравенство


$$


f(x)\overset{\text{df}}{=}{}-(1-x^2)(2x-a-b)+2x(x-a)(x-b)\ge C_9.


$$


Тогда при $x_{1,n}\ge 1-\theta_1$ имеем


\begin{equation}


\frac{1}{\sqrt{1-x_{1,n}}}\le


\frac{C_{10}}{\sqrt{(1-x^2_{1,n})(x_{1,n}-a)(x_{1,n}-b)}}=


\frac{C_{10}}{\displaystyle\int^1_{x_{1,n}}\dfrac{f(x)dx}


{2\sqrt{(1-x^2)(x-a)(x-b)}}}\le\frac{2C_{10}}{C_9\theta_2},


\end{equation}


причем $\theta_2=F(x_{1,n})$


$$


F(x)=\int^1_x\frac{dx}{\sqrt{(1-x^2)(x-a)(x-b)}},\quad x\in(b,1).


$$


Поскольку $F$ убывает на $(b,1)$ и


$S_n(F^{-1}(\theta_2))-S_n(F^{-1}(0))=\frac{\pi}{2}$, можно записать


$S'_n(F^{-1}(\theta_3))\times$ $\times\frac{\theta_2}{F'(\zeta)}=


\frac{\pi}{2}$, или $\theta_2=\frac{\pi}{2\gamma_n(\zeta)}$,


где $\zeta=F^{-1}(\theta_3)$, $x_{1,n}<\zeta<1$,


$0<\theta_3<\theta_2$.





Подставляя полученное значение в (23), имеем


\begin{equation}


\frac{1}{\sqrt{1-x_{1,n}}}\le\frac{2C_{10}\gamma_n(\zeta)}


{C_9\pi},


\end{equation}


откуда с учетом (4) получаем


\begin{multline}


\frac{\xi_k-x_{k,n}}{1-\xi_k}\le


\frac{\pi\sqrt{2(1-a)(1-b)}}{\sqrt{1-x_{1,n}}|\gamma_n(\xi_k)|}\le


C_{11}\biggl|\frac{\gamma_n(\zeta)}{\gamma_n(\xi_k)} \biggr|\le \\


%


\le\frac{C_{11}C_5}{C_4}\;


\frac{


\displaystyle\sum\limits^n_{j=1}


\sqrt{1-|a_{j,n}|}\,|\zeta-c_{j,n}|\Big/


(1-|a_{j,n}|\zeta)


}


{\displaystyle\sum\limits^n_{j=1}


\sqrt{1-|a_{j,n}|}\,|\xi_k-c_{j,n}|\Big/ (1-|a_{j,n}|\xi_k)


}\le


%


C_{12}\frac{1-a}{\xi_k-b}\le\frac{2C_{12}(1-a)}{b+1}


\end{multline}


при $\xi_k-b\ge(1-b)/2$. Если же $\xi_k-b<(1-b)/2$ или


$x_{1,n}<1-\theta_1$, то


\begin{equation}


Z_k\le\max\biggl(\sqrt{\frac{2}{1-b}},\,\frac{1}{\sqrt{\theta_1}}


\biggr).


\end{equation}





Запишем оценку


\begin{equation}


H_k\le(1+C_{13}B_k)\frac{\sqrt{x_{k,n}-b}}{\sqrt{\xi_k-b}}\cdot


\frac{\sqrt{1-a}}{\sqrt{1-b}},


\end{equation}


где


$$


B_k=\sum^n_{j=1}\frac{\sqrt{1-|a_{j,n}|}\,|x_{k,n}-\xi_k|}


{(1-|a_{j,n}|\xi_k)(1-|a_{j,n}|x_{k,n})}\biggl.\biggr/


\sum^n_{j=1}


\frac{\sqrt{1-|a_{j,n}|}\,|x_{k,n}-c_{j,n}|}


{1-|a_{j,n}|x_{k,n}}.


$$


Предположим сначала, что $x_{k,n}-b\ge\theta_4$, где $\theta_4>0$


таково, что при $b\le x\le b+\theta_4<1$


справедливо неравенство $f(x)\ge C_9$. Тогда (считая для определенности


$b>0$)


\begin{equation}


B_k\frac{\sqrt{x_{k,n}-b}}{\sqrt{\xi_k-b}}\le


\sum^{n_1}_{j=1}


\frac{\sqrt{1-|a_{j,n}|}\,|x_{k,n}-\xi_k|}


{(1-|a_{j,n}|x_{k,n})(1-|a_{j,n}|\xi_k)}\biggl.\biggr/


\biggl(\sum^n_{j=1}\frac{\sqrt{1-|a_{j,n}|}}{1-|a_{j,n}|x_{k,n}}


\cdot\frac{\theta_4}{\sqrt{1-b}}\biggr)+


\frac{\sqrt{(1-b)^3}}{\theta_4C_6}.


\end{equation}


Теперь будем оценивать величину


\begin{equation}


\Lambda_{k,j}\overset{\text{df}}{=}{}


\frac{|\xi_k-x_{k,n}|}{1-|a_{j,n}|\xi_k}\le


\frac{x_{k-1,n}-x_{k,n}}{1-|a_{j,n}|\xi_k}=


\frac{\pi\sqrt{(1-\xi^2_k)(\xi_k-a)(\xi_k-b)}}


{(1-|a_{j,n}|\xi_k)|\gamma_n(\xi_k)|}.


\end{equation}


При $0<1-\xi_k<1-\max\limits_{1\le\lambda\le n_1}|a_{\lambda,n}|<


\min((1-b)/2,\,1/2)$ из (28)


получаем


\begin{equation}


|\Lambda_{k,j}|\le


\frac{\pi\sqrt{2}\sqrt{(1-a)(1+b)}\sqrt{1-\xi_k}}


{(1-\max\limits_{1\le\lambda\le n_1}|a_{\lambda,n}|)


|\gamma_n(\xi_k)|}\le C_{14}\biggl.\biggr/


\sum^{n_1}_{j=1}


\frac{\sqrt{1-|a_{j,n}|}\,


\sqrt{1-\max\limits_{1\le\lambda\le n_1}|a_{\lambda,n}|}}


{1-|a_{j,n}|+1-\max\limits_{1\le\lambda\le n_1}|a_{\lambda,n}|}


\le\frac{C_{14}}{C_1}.


\end{equation}


Аналогично рассматриваются возможности неравенств


$$


1-\max_{1\le\lambda\le n_1}|a_{\lambda,n}|\le 1-\xi_k


\le\min\biggl(\frac{1-b}{2},\,\frac{1}{2} \biggr)\quad


\text{и}\quad


1-\xi_k>\min\biggl(\frac{1-b}{2},\,\frac{1}{2} \biggr).


$$





Неравенства (27)--(30) означают, что при $x_{k,n}-b\ge\theta_4$


\begin{equation}


H_k\le C_{15}.


\end{equation}


Пусть теперь $x_{k,n}-b\le\theta_4\le\xi_k-b$. Тогда в силу (7), (16)


\begin{multline*}


\frac{\sqrt{(x_{k,n}-b)(x_{k,n}-a)}}{\sqrt{\xi_k-b}}\cdot


\sum\limits^n_{j=1}


\frac{\sqrt{1-|a_{j,n}|}\,|x_{k,n}-\xi_k|}


{(1-|a_{j,n}|x_{k,n})|1-a_{j,n}\xi_k|}\biggl.\biggr/


\sum\limits^n_{j=1}


\frac{\sqrt{1-|a_{j,n}|}\,|x_{k,n}-c_{j,n}|}


{1-|a_{j,n}|x_{k,n}} \le \\


\le \biggl(\sum\limits^{n_1}_{j=1}


\frac{\sqrt{1-|a_{j,n}|}\,\Lambda_{k,j}}{1-(b+\theta_4)}


+\sum\limits^n_{j=n_1+1}


\frac{\sqrt{2}\sqrt{1-|a_{j,n}|}}{1-(b+\theta_4)}\biggr)\biggl.\biggr/


\biggl(C_3\sum\limits^n_{j=1}\sqrt{1-|a_{j,n}|}


\frac{b-1}{2}\sqrt{\theta_4}\biggr),


\end{multline*}


откуда с учетом (27)--(30) снова получаем (31). Аналогично разбираются


остальные варианты.





Итак, из (17)--(25) и (31) при $i=1,2,\dotsc,q$; $k=i+1,i+2,\dotsc,m$


имеем


\begin{equation}


|l_{k,n}(\frak M,x)|\le C_{16}


\frac{x_{k-1,n}-x_{k,n}}{x-x_{k,n}}.


\end{equation}


Подобным образом найдем, что (32) справедливо и при $i=1,2,\dotsc,q$;


$k=m+2,\dotsc,n$; $k=i=2,3,\dotsc,q+1$, кроме того, для


$i=3,\dotsc,q+1$; $k=1,2,\dotsc,i-2$


$$


|l_{k,n}(\frak M,x)|\le\frac{C_{17}}{i-k-1},


$$


а при $k=m+1$


$$


|l_{k,n}(\frak M,x)|\le C_{18}.


$$


Из полученных неравенств с учетом соотношения


$|l_{1,n}(\frak M,x)|\le 1+\sum\limits^n_{k=2}


|l_{k,n}(\frak M,x)|$ найдем оценку


$$


\sum^n_{k=1}|l_{k,n}(\frak M,x)|\le C_{19}\ln(i+1)+


2C_{16}\sum^n\begin{Sb}k=i+2 \\ k\ne m+1\end{Sb}


\frac{x_{k-1,n}-x_{k,n}}{x-x_{k,n}}.


$$


Для завершения доказательства теоремы осталось заметить, что


(см. также (24))


\begin{multline*}


\sum^n\begin{Sb}k=i+2 \\ k\ne m+1\end{Sb}


\frac{x_{k-1,n}-x_{k,n}}{x-x_{k,n}}=


\sum^n\begin{Sb}k=i+2 \\ k\ne m+1\end{Sb}


\int^{x_{k-1,n}}_{x_{k,n}}\frac{d\theta}{x-x_{k,n}}\le


\sum^n\begin{Sb}k=i+2 \\ k\ne m+1\end{Sb}


\int^{x_{k-1,n}}_{x_{k,n}}\frac{d\theta}{x-\theta}\le


\int^{x_{i+1,n}}_{x_{n,n}}\frac{d\theta}{x-\theta}-\\


-\int^b_{x_{m+1,n}}\frac{d\theta}{x-\theta}\le


\ln\frac{2}{x_{i,n}-x_{i+1,n}}=\ln\,


\frac{2\gamma_n(\xi_{i+1})}


{\sqrt{(1-\xi^2_{i+1})(\xi_{i+1}-a)(\xi_{i+1}-b)}}


\le\ln\|\gamma_n\|^2_{C(E)}.


\end{multline*}





Автор глубоко благодарен своему покойному учителю, профессору


А.А.\,Привалову за постановку задачи.





\begin{thebibliography}{99}





\bibitem{1}


Привалов А.А. {\em Теория интерполирования функций}. Кн.\,1. --


Саратов: Изд-во СГУ, 1990. -- 230~с.


\bibitem{2}


Турецкий А.Х. {\em Теория интерполирования в задачах}. -- Минск:


Вышэйш. школа, 1968. -- 318~с.


\bibitem{3}


Ибрагимов И.И. {\em Методы интерполяции функций и некоторые их


применения}. -- М.: Наука, 1971. -- 518~с.


\bibitem{4}


Уолш Дж.Л.{\em Интерполяция и аппроксимация рациональными функциями в


комплексной области}. -- М.: Ин. лит., 1961. -- 508~с.


\bibitem{5}


Русак В.Н. {\em Об интерполировании рациональными функциями с


фиксированными полюсами} // ДАН БССР. -- 1962. -- Т.\,6. -- \No\,9. --


С.\,548--550.


\bibitem{6}


Ровба Е.А. {\em Приближение выпуклых функций интерполяционными


рациональными функциями с фиксированным числом полюсов} // ДАН БССР. --


1977. -- Т.\,21. -- \No\,9. -- С.\,781--783.


\bibitem{7}


Ровба Е.А. {\em О рациональной интерполяции функции $|x|$} // Изв. АН


БССР. Сер. физ.-матем. наук. -- 1989. -- \No\,5. -- С.\,39--46.


\bibitem{8}


Старовойтов А.П. {\em О рациональной интерполяции с фиксированными


полюсами}. -- Ред. журн. ``Изв. АН БССР. Сер. физ.-матем. наук.'' --


Минск, 1983. -- 18~с. -- Деп. в ВИНИТИ 22.05.83, \No\,2735-83.


\bibitem{9}


Лукашов А.Л. {\em О задаче Чебышева-Маркова на двух отрезках}. --


Сарат. ун-т. -- Саратов, 1989. -- 56~с. -- Деп. в ВИНИТИ 01.11.89,


\No\,6615-B89.


\bibitem{10}


Аптекарев А.И. {\em Асимптотические свойства многочленов, ортогональных


на системе контуров, и периодические движения цепочек Тода} // Матем.


сб. -- 1984. -- Т.\,125. -- \No\,2. -- С.\,231--258.


\bibitem{11}


Peherstorfer F. {\em Orthogonal and Chebyshev polynomials on two


intervals} // Acta Math. Hung. -- 1990. -- V.\,55. -- \No\,3--4. --


P.\,245--278.


\bibitem{12}


Peherstorfer F. {\em Elliptic orthogonal and extremal polynomials} //


Proc. London Math. Soc. -- 1995. -- V.\,70. -- P.\,605--624.


\bibitem{13}


Содин М.Л., Юдицкий П.М. {\em Алгебраическое решение задач


Е.И.\,Золотарева и Н.И.\,Ахиезера о многочленах, наименее уклоняющихся


от нуля} // Теория функций, функц. анализ и их прилож. -- Харьков,


1991. -- Вып.\,56. -- С.\,56--64.


\bibitem{14}


Бернштейн С.Н. {\em Собрание сочинений}. Т.\,1. {\em Конструктивная


теория функций}. -- М.: Изд-во АН СССР, 1952. -- 582~с.


\bibitem{15}


Kilgore T.A. {\em A note on optimal interpolation with rational


functions} // Acta sci. math. -- 1988. -- V.\,52. -- P.\,113--116.


\bibitem{16}


Привалов И.И. {\em Введение в теорию функций комплексного переменного}.


-- 12-е изд. -- М.: Наука, 1977. -- 444~с.








\end{thebibliography}





\vskip 0.5cm





\post{\it Саратовский государственный}{\it Поступила}


\post{\it университет}{25.07.1997}





\label{end}


\end{document}


